
The validation attempt produced one artifact---HumanEval + Pydantic type extensions---while all hypothesis testing was blocked by infrastructure failures.

\subsubsection{Dataset Extension: Artifact Contribution}

The dataset extension pipeline successfully generated 100 Pydantic-annotated prompts from HumanEval problems 0--99.

\begin{table*}[htbp]
\centering
\resizebox{\dimexpr 2\columnwidth+\columnsep\relax}{!}{%
\begin{tabular}{lllll}
\toprule
Component & Input & Expected Output & Actual Output & Status \\
\midrule
DatasetExtender & HumanEval (100 problems) & 100 Pydantic prompts & 100 JSON files & ✓ PASS \\
\bottomrule
\end{tabular}
}
\end{table*}

Manual inspection of 10 randomly-sampled prompts confirmed type annotations match original signatures, validators encode docstring preconditions, and test cases remain unchanged.

\textbf{Artifact Availability}: Dataset available in \texttt{src/data/humaneval\textbackslash \_pydantic/}. This artifact demonstrates that typed benchmarks can be mechanically generated from existing code generation datasets.

\subsubsection{LLM Generation: Infrastructure Failure}

All LLM code generation attempts failed with HTTP 401 authentication errors.

\textbf{Quantitative Results}:
\begin{itemize}
\item Total generation attempts: 48
\item Successful generations: 0 (0\% success rate)
\item Error type: HTTP 401 ''API key is invalid'' (100\% of failures)
\item Retry attempts per request: 3 (exponential backoff: 1s, 2s, 4s)
\end{itemize}

\textbf{Error Message}:
\begin{verbatim}
Error code: 401 - {'type': 'error', 'error': {'type': 'authentication_error', 
'message': 'API key is invalid.'}}
\end{verbatim}

\textbf{Root Cause Analysis}: Three competing explanations were evaluated. (1) Invalid API key: Error message explicitly states ''API key is invalid''; consistent across all 48 attempts. Plausibility: HIGH. (2) API rate limiting: Would produce HTTP 429 errors, not 401. Plausibility: LOW. (3) Network/firewall blocking: Would produce connection timeouts, not authenticated 401 responses. Plausibility: LOW.

\textbf{Most Likely Interpretation}: Invalid API key. Retry with confirmed-valid key would distinguish between infrastructure failure and hypothesis refutation.

The h-e1 gate criterion (extraction\_success\_rate $\geq 90$\%) could not be tested. The hypothesis remains theoretically plausible but empirically unvalidated. All downstream hypotheses (h-m1, h-m2, h-m3) were blocked.

\subsubsection{Pipeline Architectural Validation}

Despite generation failure, the constraint extraction pipeline was validated on error files.

\begin{table*}[htbp]
\centering
\resizebox{\dimexpr 2\columnwidth+\columnsep\relax}{!}{%
\begin{tabular}{lllll}
\toprule
Component & Input & Expected Output & Actual Output & Status \\
\midrule
DatasetExtender & HumanEval (100 problems) & 100 Pydantic prompts & 100 JSON files & ✓ PASS \\
LLMGenerator & 100 Pydantic prompts & 100 typed programs & 48 error files & ✗ FAIL (infrastructure) \\
PyreExtractor & 48 error files & 0 constraints & 0 constraints & ✓ PASS (negative case) \\
Z3Validator & 0 constraints & 0 quality constraints & 0 quality constraints & ✓ PASS (negative case) \\
MetricsEngine & 0 constraints & extraction\_rate=0\%, gate FAIL & extraction\_rate=0.0\%, gate FAIL & ✓ PASS \\
\bottomrule
\end{tabular}
}
\end{table*}

\textbf{Architectural Soundness}: Pipeline validated on negative cases only. PyreExtractor AST parser correctly handles unparseable Python, returns 0 constraints as expected. Z3Validator correctly identifies 0 quality constraints from empty input. MetricsEngine correctly computes 0\% extraction rate and triggers FAIL condition.

The pipeline is architecturally sound and ready for retry. Fixing the API key requires no code changes---only environment configuration.

\subsubsection{Visualization Results}

Four figures were generated, though three contain placeholder/zero data.

\textbf{Figure 1 (gate\_metrics.png)}: Gate threshold (90.0\%) vs actual extraction rate (0.0\%). Shows gate failure.

\textbf{Figure 2 (failure\_modes.png)}: Categorization of extraction failures. 100\% API authentication errors (48/48 attempts). This is the only figure with meaningful data.

\textbf{Figure 3 (success\_by\_complexity.png)}: Planned analysis of extraction success by program complexity. All bins show 0\% due to no successful generations.

Figure 2 is the critical evidence: the failure mode is homogeneous (100\% authentication errors), not heterogeneous. This supports the interpretation that infrastructure failure, not scientific invalidity, blocked validation.

\subsubsection{Hypothesis Validation Summary}

\begin{table*}[htbp]
\centering
\resizebox{\dimexpr 2\columnwidth+\columnsep\relax}{!}{%
\begin{tabular}{llllll}
\toprule
Hypothesis & Gate & Threshold & Actual Result & Status & Reason \\
\midrule
h-e1 & MUST\_WORK & extraction\_rate $\geq 90$\% & 0.0\% & ✗ FAIL & Infrastructure block (API auth) \\
h-m1 & MUST\_WORK & All programs extract constraints & N/A & NOT\_STARTED & Prerequisite h-e1 failed \\
h-m2 & SHOULD\_WORK & cone\_size <30\%, detection\_rate=100\% & N/A & NOT\_STARTED & Prerequisite h-e1, h-m1 failed \\
h-m3 & MUST\_WORK & median speedup $\geq 2x$ & N/A & NOT\_STARTED & Prerequisite h-e1, h-m1, h-m2 failed \\
\bottomrule
\end{tabular}
}
\end{table*}

\textbf{Assumptions}:
\begin{itemize}
\item A1 (LLM code has extractable constraints): UNVERIFIED (h-e1 blocked)
\item A2 (repair locality): UNVERIFIED (h-m2 not reached)
\item A3 (dependency analysis soundness): UNVERIFIED (h-m2 not reached)
\item A4 (extraction overhead small): UNVERIFIED (h-m1 not reached)
\item A5 (no dynamic features): PARTIALLY VERIFIED (prompt constraints applied, not tested on actual LLM output)
\end{itemize}
